% Problem record op_9aba71db480b451a. This TeX file is derived from
% database/problems_json/op_9aba71db480b451a.json by scripts/sync-tex.mjs; edit the JSON record
% and rerun the script rather than editing this file.
\section{Pseudorandom unitaries from ordinary random two-qubit circuits}
\label{sec:op_9aba71db480b451a}

\paragraph{Problem.}
Does a polynomial-length random walk generated by ordinary two-qubit gates form a pseudorandom-unitary ensemble?

For an ensemble $\mathcal U_n$ of efficiently implementable $n$-qubit unitaries, require that every quantum polynomial-time distinguisher with forward oracle access satisfy

\begin{equation}
\left|\Pr_{U\sim\mathcal U_n}[\mathcal D^U(1^n)=1]
-\Pr_{V\sim\operatorname{Haar}(2^n)}[\mathcal D^V(1^n)=1]\right|
\leq\operatorname{negl}(n).
\label{eq:9aba-1}
\end{equation}

Here $\operatorname{negl}(n)<n^{-a}$ eventually for every constant $a>0$. One hidden unitary is sampled and reused; inverse, transpose, conjugate, and controlled-unitary oracles are not supplied.

Fix a finite universal gate set $\mathcal G\subset SU(4)$. At each step choose a qubit pair and a gate from $\mathcal G$ independently and uniformly, and let $U_{n,m}$ be the resulting circuit. Do constants $c,C_{\mathcal G}>0$ exist such that

\begin{equation}
\mathcal U_n=\{U_{n,m(n)}:m(n)=\lceil C_{\mathcal G}n^c\rceil\}
\label{eq:9aba-2}
\end{equation}

satisfies the pseudorandomness condition in Eq.~\eqref{eq:9aba-1}? Equation~\eqref{eq:9aba-2} fixes one polynomial-size ensemble against all polynomial-time distinguishers.

\subsection*{Status}

Unsolved

\subsection*{Source}

The question is explicitly posed or retained as open in the cited primary literature \sourcecite{ref:9aba-ji18}{Ji18}\sourcecite{ref:9aba-bostanci25}{Bostanci25}\sourcecite{ref:9aba-raza26}{Raza26}. The statement is rewritten here to make its hypotheses and success criterion self-contained.

\subsection*{Progress}

\begin{itemize}
  \item Status: source-explicit conjecture, with a September 2026 consistency check. The HPS paper states this conjecture rather than proving it. The introduction of the September 2026 distinctness paper still describes cryptographic pseudorandomness of random circuits as a belief, not an established general theorem. \sourcecite{ref:9aba-bostanci25}{Bostanci25}\sourcecite{ref:9aba-raza26}{Raza26}

  \item There are strong positive results for approximate designs and for specially constructed low-depth PRUs. Schuster, Haferkamp, and Huang obtain the latter by assembling patches drawn from existing PRU constructions. This does not identify their ensemble with the independent, uniformly sampled two-qubit-gate walk above. \sourcecite{ref:9aba-schuster25}{Schuster25}

  \item No subsequent proof or efficient attack resolving this exact random-walk conjecture was located in the literature check.
\end{itemize}

\subsection*{References}

\begin{enumerate}[label={},leftmargin=4.8em,labelsep=0.5em]
  \footnotesize
  \item[\textup{[Ji18]}]\label{ref:9aba-ji18}
  Zhengfeng Ji, Yi-Kai Liu, and Fang Song, \emph{Pseudorandom Quantum States}. \href{https://eprint.iacr.org/2018/544.pdf}{Cryptology ePrint 2018/544}, CRYPTO 2018. See §6.1 for the PRU definition and §6.2, especially printed p. 23, for the constant-round candidate with different keys.

  \item[\textup{[Bostanci25]}]\label{ref:9aba-bostanci25}
  Bostanci, Haferkamp, Hangleiter, and Poremba, \emph{Efficient Quantum Pseudorandomness from Hamiltonian Phase States}. \href{https://arxiv.org/html/2410.08073}{arXiv:2410.08073}, Conjecture 1.1.

  \item[\textup{[Raza26]}]\label{ref:9aba-raza26}
  Raza, Eisert, and Fefferman, \href{https://arxiv.org/html/2609.03065v1}{arXiv:2609.03065v1}, Introduction.

  \item[\textup{[Schuster25]}]\label{ref:9aba-schuster25}
  Thomas Schuster, Jonas Haferkamp, and Hsin-Yuan Huang, \emph{Random unitaries in extremely low depth}. \href{https://arxiv.org/html/2407.07754}{arXiv:2407.07754}; Science 389, 92 (2025). See the construction framework and Corollary 2 for the cryptographic assumptions and the use of local PRU ingredients.
\end{enumerate}

\subsection*{Comment}

The mathematical bottleneck is the order of quantifiers. Proving that circuits of size $m(n,t)$ emulate $t$ Haar moments does not show that one fixed polynomial-size ensemble defeats every polynomial-time observer. For example, establishing a bound only through $t=n^a$ leaves adversaries using $n^{a+1}$ queries outside the theorem.

Similarly, hardness of classically sampling an output distribution would not by itself establish indistinguishability against a quantum algorithm allowed coherent, adaptive access to the operation.

The question asks whether generic random evolution has cryptographic security, rather than whether a specially engineered cryptographic circuit can imitate random evolution. It is distinct from structured phase--Hadamard pseudorandom-unitary constructions.

\subsection*{Field}

Quantum Cryptography; Quantum algorithm

\subsection*{Topic}

Quantum circuit complexity; Computational complexity and computability

\subsection*{ID}

\texttt{op\_9aba71db480b451a}
